% GENERATED FILE. Edit canonical lessons, then run npm run generate:books.
% canonical-source-hash: fnv1a64:ff57000bb465938c
% canonical-lessons: HI-C520-sulgana, HI-C520-gudhna, HI-C520-cakla, HI-C520-belna, HI-C520-phulna

\chapter{Making Rotis}
\label{ch:hi-making-rotis}
\hlchaptermodality{520}

\begin{chapteropening}
\textbf{By the end of this chapter:} \emph{I can describe lighting the stove and making rotis.}

Complete the last lesson of chapter 520: I can describe lighting the stove and making rotis.
\end{chapteropening}

\section[\texorpdfstring{\hi{सुलगाना} (sulgānā)}{सुलगाना (sulgānā)}]{\hi{सुलगाना} (sulgānā) — to kindle; to light (a fire)}

\label{lesson:HI-C520-sulgana}



\begin{quote}
Before the new one: say the Hindi for to sweep (with a broom), then the Hindi for to scour (pots).
\end{quote}

\subsection*{You'll want to know: \hi{सुलगाना}}
\textbf{\hi{सुलगाना}} — \emph{sulgānā} — \textquotedblleft{}to kindle; to light (a fire)\textquotedblright{}.

\begin{quote}
\hi{दादी} \hi{ने} \hi{सुबह} \hi{चूल्हा} \hi{सुलगाया।}
\end{quote}

\begin{quote}
\emph{Dādī ne subah cūlhā sulgāyā.}
\end{quote}

\begin{quote}
Grandmother lit the stove in the morning.
\end{quote}

Past, with \hi{ने} (he or she …): \hi{उसने} \hi{आग} \hi{सुलगाई}, \emph{usne āg sulgāī}.

\begin{cousinweb}
\textbf{\hi{सुलगाना}} (\emph{sulgānā}) is the causative of \emph{sulagnā}, \textquotedblleft{}to catch fire\textquotedblright{}, which CDIAL files under Sanskrit \emph{saṁlagna}, \textquotedblleft{}attached, joined\textquotedblright{} (CDIAL 12999).
\end{cousinweb}

\begin{grammarlens}[title={What you've built}]
One more word for describing what you lived and felt.
\end{grammarlens}

\subsection*{Your turn}
\begin{itemize}
\raggedright
  \item \emph{Say it:} \emph{sulgānā}
  \item \emph{Say it:} \emph{sulgānā}, once more
  \item \emph{Say it:} \emph{Dādī ne subah cūlhā sulgāyā.} — Grandmother lit the stove in the morning.
  \item \emph{Recall:} say the Hindi for to sweep (with a broom), then the Hindi for to scour (pots), then say \emph{sulgānā} again
\end{itemize}

\begin{tcolorbox}[breakable,colback=teal!4,colframe=teal!35!black,title={Before you move on}]
What does \hi{सुलगाना} mean? (\textquotedblleft{}To kindle; to light (a fire)\textquotedblright{}.) Say it once more.
\end{tcolorbox}
\section[\texorpdfstring{\hi{गूँधना} (gū̃dhnā)}{गूँधना (gū̃dhnā)}]{\hi{गूँधना} (gū̃dhnā) — to knead}

\label{lesson:HI-C520-gudhna}



\begin{quote}
Before the new one: say the Hindi for to scour (pots), then the Hindi for to kindle.
\end{quote}

\subsection*{You'll want to know: \hi{गूँधना}}
\textbf{\hi{गूँधना}} — \emph{gū̃dhnā} — \textquotedblleft{}to knead\textquotedblright{}.

\begin{quote}
\hi{माँ} \hi{आटा} \hi{गूँध} \hi{रही} \hi{है।}
\end{quote}

\begin{quote}
\emph{Māṁ āṭā gū̃dh rahī hai.}
\end{quote}

\begin{quote}
Mother is kneading the dough.
\end{quote}

Past, with \hi{ने} (he or she …): \hi{उसने} \hi{आटा} \hi{गूँधा}, \emph{usne āṭā gū̃dhā}.

\begin{cousinweb}
\textbf{\hi{गूँधना}} (\emph{gū̃dhnā}) is a tadbhava, inherited from Sanskrit \emph{guphati}, \textquotedblleft{}strings together\textquotedblright{} (CDIAL 4205). CDIAL gives the Hindi verb as \textquotedblleft{}to plait, mix, knead\textquotedblright{}.
\end{cousinweb}

\begin{grammarlens}[title={What you've built}]
One more word for describing what you lived and felt.
\end{grammarlens}

\subsection*{Your turn}
\begin{itemize}
\raggedright
  \item \emph{Say it:} \emph{gū̃dhnā}
  \item \emph{Say it:} \emph{gū̃dhnā}, once more
  \item \emph{Say it:} \emph{Māṁ āṭā gū̃dh rahī hai.} — Mother is kneading the dough.
  \item \emph{Recall:} say the Hindi for to scour (pots), then the Hindi for to kindle, then say \emph{gū̃dhnā} again
\end{itemize}

\begin{tcolorbox}[breakable,colback=teal!4,colframe=teal!35!black,title={Before you move on}]
What does \hi{गूँधना} mean? (\textquotedblleft{}To knead\textquotedblright{}.) Say it once more.
\end{tcolorbox}
\section[\texorpdfstring{\hi{चकला} (caklā)}{चकला (caklā)}]{\hi{चकला} (caklā) — a round board for rolling dough}

\label{lesson:HI-C520-cakla}



\begin{quote}
Before the new one: say the Hindi for to kindle, then the Hindi for to knead.
\end{quote}

\subsection*{You'll want to know: \hi{चकला}}
\textbf{\hi{चकला}} — \emph{caklā} — \textquotedblleft{}a round board for rolling dough\textquotedblright{}.

\begin{quote}
\hi{चकला} \hi{और} \hi{बेलन} \hi{मेज़} \hi{पर} \hi{रखे} \hi{हैं।}
\end{quote}

\begin{quote}
\emph{Caklā aur belan mez par rakhe haiṁ.}
\end{quote}

\begin{quote}
The rolling board and the rolling pin are on the table.
\end{quote}

Masculine.

\begin{cousinweb}
\textbf{\hi{चकला}} (\emph{caklā}) is a tadbhava, inherited from \emph{*cakrala}, \textquotedblleft{}round and flat\textquotedblright{}, a form CDIAL reconstructs from its descendants (CDIAL 4548).
\end{cousinweb}

\begin{grammarlens}[title={What you've built}]
One more word for describing what you lived and felt.
\end{grammarlens}

\subsection*{Your turn}
\begin{itemize}
\raggedright
  \item \emph{Say it:} \emph{caklā}
  \item \emph{Say it:} \emph{caklā}, once more
  \item \emph{Say it:} \emph{Caklā aur belan mez par rakhe haiṁ.} — The rolling board and the rolling pin are on the table.
  \item \emph{Recall:} say the Hindi for to kindle, then the Hindi for to knead, then say \emph{caklā} again
\end{itemize}

\begin{tcolorbox}[breakable,colback=teal!4,colframe=teal!35!black,title={Before you move on}]
What does \hi{चकला} mean? (\textquotedblleft{}A round board for rolling dough\textquotedblright{}.) Say it once more.
\end{tcolorbox}
\section[\texorpdfstring{\hi{बेलना} (belnā)}{बेलना (belnā)}]{\hi{बेलना} (belnā) — to roll out (dough)}

\label{lesson:HI-C520-belna}



\begin{quote}
Before the new one: say the Hindi for to knead, then the Hindi for a round board for rolling dough.
\end{quote}

\subsection*{You'll want to know: \hi{बेलना}}
\textbf{\hi{बेलना}} — \emph{belnā} — \textquotedblleft{}to roll out (dough)\textquotedblright{}.

\begin{quote}
\hi{मीरा} \hi{ने} \hi{चकले} \hi{पर} \hi{गोल} \hi{रोटी} \hi{बेली।}
\end{quote}

\begin{quote}
\emph{Mīrā ne cakle par gol roṭī belī.}
\end{quote}

\begin{quote}
Mira rolled out a round roti on the board.
\end{quote}

Past, with \hi{ने} (he or she …): \hi{उसने} \hi{पराठा} \hi{बेला}, \emph{usne parāṭhā belā}.

\begin{cousinweb}
\textbf{\hi{बेलना}} (\emph{belnā}) is a tadbhava, inherited from Sanskrit \emph{vēllayati}, \textquotedblleft{}causes to shake, kneads\textquotedblright{} (CDIAL 12121). CDIAL gives the Hindi verb as \textquotedblleft{}to roll out\textquotedblright{}.
\end{cousinweb}

\begin{grammarlens}[title={What you've built}]
One more word for describing what you lived and felt.
\end{grammarlens}

\subsection*{Your turn}
\begin{itemize}
\raggedright
  \item \emph{Say it:} \emph{belnā}
  \item \emph{Say it:} \emph{belnā}, once more
  \item \emph{Say it:} \emph{Mīrā ne cakle par gol roṭī belī.} — Mira rolled out a round roti on the board.
  \item \emph{Recall:} say the Hindi for to knead, then the Hindi for a round board for rolling dough, then say \emph{belnā} again
\end{itemize}

\begin{tcolorbox}[breakable,colback=teal!4,colframe=teal!35!black,title={Before you move on}]
What does \hi{बेलना} mean? (\textquotedblleft{}To roll out (dough)\textquotedblright{}.) Say it once more.
\end{tcolorbox}
\section[\texorpdfstring{\hi{फूलना} (phūlnā)}{फूलना (phūlnā)}]{\hi{फूलना} (phūlnā) — to puff up; to swell; to bloom}

\label{lesson:HI-C520-phulna}



\begin{quote}
Before the new one: say the Hindi for a round board for rolling dough, then the Hindi for to roll out (dough).
\end{quote}

\subsection*{You'll want to know: \hi{फूलना}}
\textbf{\hi{फूलना}} — \emph{phūlnā} — \textquotedblleft{}to puff up; to swell; to bloom\textquotedblright{}.

\begin{quote}
\hi{गरम} \hi{तवे} \hi{पर} \hi{रोटी} \hi{फूल} \hi{गई।}
\end{quote}

\begin{quote}
\emph{Garam tave par roṭī phūl gaī.}
\end{quote}

\begin{quote}
The roti puffed up on the hot griddle.
\end{quote}

Past (it …): \hi{वह} \hi{रोटी} \hi{फूल} \hi{गई}, \emph{vah roṭī phūl gaī}.

\begin{cousinweb}
\textbf{\hi{फूलना}} (\emph{phūlnā}) is a tadbhava, inherited from Sanskrit \emph{phullati}, \textquotedblleft{}expands, opens (as a flower)\textquotedblright{} (CDIAL 9093). CDIAL gives the Hindi verb as \textquotedblleft{}to swell, blossom\textquotedblright{}.
\end{cousinweb}

\begin{grammarlens}[title={What you've built}]
One more word for describing what you lived and felt.
\end{grammarlens}

\subsection*{Your turn}
\begin{itemize}
\raggedright
  \item \emph{Say it:} \emph{phūlnā}
  \item \emph{Say it:} \emph{phūlnā}, once more
  \item \emph{Say it:} \emph{Garam tave par roṭī phūl gaī.} — The roti puffed up on the hot griddle.
  \item \emph{Recall:} say the Hindi for a round board for rolling dough, then the Hindi for to roll out (dough), then say \emph{phūlnā} again
\end{itemize}

\begin{tcolorbox}[breakable,colback=teal!4,colframe=teal!35!black,title={Before you move on}]
What does \hi{फूलना} mean? (\textquotedblleft{}To puff up; to swell; to bloom\textquotedblright{}.) Say it once more.
\end{tcolorbox}
